\par\Needspace{8\baselineskip}\bigskip
\pdfbookmark[3]{Exposition, 13}{pallas-II-article13}
\begin{center}13.\end{center}
\noindent\hspace*{\parindent}Under the usual convention, the epoch of mean longitude for time \(0\) is \(L-nt\), and this would give
\[
\begin{aligned}
&\mathrm{d}(\text{Epoch of mean long.})=-t\mathrm{d}n+(1-\cos i)\mathrm{d}\AscendingNode\\
&\quad-(2ar+aa\cos\varphi\mathbin{.}\operatorname{tang}\tfrac12\varphi\mathbin{.}\cos(v-\tilde\omega))Tn\mathrm{d}t\\
&\quad+\dfrac{ar\operatorname{tang}\tfrac12\varphi\mathbin{.}\sin(v-\tilde\omega)\mathbin{.}(2+e\cos(v-\tilde\omega))}{\cos\varphi}Vn\mathrm{d}t,
\end{aligned}
\]
where for\raisebox{-.45ex}{-} \(\mathrm{d}n\) one could substitute the value found above. In this form, however, the instantaneous variation in the epoch would contain a term that could grow without bound because of the factor \(t\). To avoid this difficulty, we shall define the epoch of mean longitude by \(L-\int n\mathrm{d}t\), which, as is clear, agrees with the usual convention for elliptic motion. Denoting the epoch of mean longitude for time zero by \(\epsilon\), we then have
\[
\begin{array}{rl}
70)\quad&\begin{aligned}[t]\mathrm{d}\epsilon&=(1-\cos i)\mathrm{d}\AscendingNode\\
&\quad-(2ar+aa\cos\varphi\mathbin{.}\operatorname{tang}\tfrac12\varphi\mathbin{.}\cos(v-\tilde\omega))Tn\mathrm{d}t\\
&\quad+\dfrac{ar\operatorname{tang}\tfrac12\varphi\mathbin{.}\sin(v-\tilde\omega)\mathbin{.}(2+e\cos(v-\tilde\omega))}{\cos\varphi}Vn\mathrm{d}t,\end{aligned}
\end{array}
\]
and the mean longitude for time \(t\) will be determined by the formula
\[
\begin{array}{rl}
71)\quad&L=\epsilon+\int n\mathrm{d}t.
\end{array}
\]
